\section{First moment versus all moments: a self-contained derivation}
\label{app:moments}
% Provenance: every identity and bound below is checked on seeded random instances by
% ../../experiments/33_moment_derivation_checks_what_correlation_and_ce_identify_start_25Nov2026/verify_moment_derivation.py.
% Its printed output is results/verify_moment_derivation_output.txt (copied into assets_log.txt by make_memo_assets.py).

This appendix shows why a correlation objective constrains one conditional moment of the human rating distribution, while cross-entropy (CE) constrains all of them.
\Cref{app:proofs} proves the identification statements of Section~\ref{sec:identification}.
Here we add five derivations: a tower-property factorization of the correlation with individual ratings, a bound that converts the CE gap into an error bound for every bounded moment, the equivalence of ``all moments'' with ``the distribution'', the logit gradient of every objective, and an exactly solvable finite-sample example.
Experiment~33 checks each identity and bound on seeded random instances with exact finite sums, and each subsection reports the result of its check.
No rating is sampled and no judge is fitted.

\subsection{Notation}
\label{app:moments-notation}

\begin{table}[h]
\caption{\textbf{Every quantity in this appendix is a function of the conditional rating distribution.}
$K$ is the number of rating categories; the numerical checks use $K=6$.}
\label{tab:notation}
\centering
\small
\begin{tabular}{@{}ll@{}}
\toprule
Symbol & Meaning \\
\midrule
$X$, $R\in\{0,\dots,K-1\}$ & an item and one individual human rating of it \\
$p(\cdot\mid x)$, $q(\cdot\mid x)$ & true conditional rating distribution and the judge's forecast \\
$\mu(x)$, $m(x)$ & true conditional mean $\E[R\mid x]$ and forecast mean $\E_q[R\mid x]$ \\
$\sigma^2(x)$ & true conditional variance $\Var(R\mid x)$ \\
$f$, $\Delta_f$ & a function of the rating and its width $\max_rf(r)-\min_rf(r)$ \\
$m_f(q,x)$ & forecast expectation $\E_q[f(R)\mid x]$; $f(r)=(r/(K-1))^j$ gives the $j$th moment \\
$\eta$ & correlation ratio $\sqrt{\Var\mu(X)/\Var R}$ \\
$\CE(q)$ & $\E[-\log q(R\mid X)]$ in nats \\
$\TVD(p,q)$ & $\tfrac12\sum_r|p(r)-q(r)|$ \\
\bottomrule
\end{tabular}
\end{table}

\subsection{Correlation sees one moment}
\label{app:moments-tower}

Section~\ref{sec:identification} correlates forecast means with item means.
A judge is also scored by its correlation with individual ratings, one rating per data point.
Both versions factor through the conditional mean.

\begin{proposition}[Correlation with individual ratings factors through the mean]
\label{prop:tower}
Let $m(X)$ be any forecast mean with $\Var m(X)>0$, and let $\eta=\sqrt{\Var\mu(X)/\Var R}$.
Then
\begin{equation}
\operatorname{Corr}\bigl(m(X),R\bigr)=\operatorname{Corr}\bigl(m(X),\mu(X)\bigr)\,\eta,
\qquad
\eta^2=1-\frac{\E\,\sigma^2(X)}{\Var R}.
\label{eq:tower}
\end{equation}
The maximum over forecast means equals $\eta$.
It is attained exactly when $m(X)=a+b\,\mu(X)$ with $b>0$, and it lies below one whenever $\E\,\sigma^2(X)>0$.
\end{proposition}

\begin{proof}
Iterated expectation gives $\E[m(X)R]=\E[m(X)\mu(X)]$ and $\E R=\E\mu(X)$, so $\Cov(m(X),R)=\Cov(m(X),\mu(X))$.
The law of total variance gives $\Var R=\Var\mu(X)+\E\sigma^2(X)$, which yields the second form of $\eta^2$.
Divide the covariance by $\operatorname{sd}(m(X))\operatorname{sd}(R)$ and write $\operatorname{sd}(R)=\operatorname{sd}(\mu(X))/\eta$.
Cauchy--Schwarz bounds $\operatorname{Corr}(m(X),\mu(X))$ by one, with equality exactly for positive affine $m$.
\end{proof}

Only $\mu$ enters Equation~\eqref{eq:tower}.
Two forecasts with the same mean function have the same correlation however their variance, skewness or tails differ.
Even the true distribution $p$ scores $\eta$, not one.
The rating noise $\E\sigma^2(X)$ caps the correlation any judge can reach, and the objective never asks the judge to fit it.

\paragraph{Check.}
On 200 random instances (seven items with Dirichlet-drawn rating distributions, $K=6$), the factorization, the law of total variance and the affine ceiling each held to within $10^{-13}$.
No random forecast exceeded the ceiling, which ranged from $0.092$ to $0.681$ (mean $0.400$) across the instances.

\subsection{Cross-entropy controls every bounded moment}
\label{app:moments-ce}

\begin{lemma}[The CE gap is an expected KL divergence]
\label{lem:cekl}
$\CE(q)-\CE(p)=\E_X\KL\bigl(p(\cdot\mid X)\,\|\,q(\cdot\mid X)\bigr)\ge0$.
\end{lemma}

\begin{proof}
Write $-\log q=-\log p+\log(p/q)$ and take the expectation under the joint law of $(X,R)$.
\end{proof}

\begin{proposition}[CE bounds the error of every bounded functional]
\label{prop:bounded}
For every $f:\{0,\dots,K-1\}\to\mathbb R$ with width $\Delta_f$,
\begin{equation}
\E_X\bigl[(m_f(q,X)-m_f(p,X))^2\bigr]\le\frac{\Delta_f^2}{2}\bigl(\CE(q)-\CE(p)\bigr).
\label{eq:bounded}
\end{equation}
\end{proposition}

\begin{proof}
Fix $x$ and write $p$ and $q$ for the conditional distributions.
With $c=\tfrac12(\max f+\min f)$ we have $|f(r)-c|\le\Delta_f/2$, so $|\E_pf-\E_qf|=\bigl|\sum_r(p(r)-q(r))(f(r)-c)\bigr|\le\Delta_f\,\TVD(p,q)$.
Pinsker's inequality (\Cref{app:proofs}) gives $\TVD(p,q)\le\sqrt{\KL(p\|q)/2}$.
Squaring and averaging over $X$ with \Cref{lem:cekl} proves Equation~\eqref{eq:bounded}.
\end{proof}

\begin{proposition}[Every integer moment and the variance]
\label{prop:variance}
(i) For $f(r)=(r/(K-1))^j$ the width is $\Delta_f=1$, and Equation~\eqref{eq:bounded} is Proposition~\ref{prop:moments}.
(ii) The conditional variance satisfies $|\sigma_p^2(x)-\sigma_q^2(x)|\le(K-1)^2\,\TVD(p,q)$ and
\[
\E_X\bigl[(\sigma_p^2(X)-\sigma_q^2(X))^2\bigr]\le\frac{(K-1)^4}{2}\bigl(\CE(q)-\CE(p)\bigr).
\]
\end{proposition}

\begin{proof}
Part (i) is immediate.
For part (ii), $\Var_p(R)=\min_a\E_p(R-a)^2$ with the minimum over $a\in[0,K-1]$.
Each $f_a(r)=(r-a)^2$ takes values in $[0,(K-1)^2]$, so $|\E_pf_a-\E_qf_a|\le(K-1)^2\TVD(p,q)$ for every $a$, and the two minima differ by at most this amount.
Squaring, applying Pinsker's inequality and averaging over $X$ finish the proof.
\end{proof}

The bounds run in one direction: a small CE gap forces small moment errors, with no assumption on the forecast beyond positivity.
The converse fails, because a forecast can match the mean and the variance and still differ in shape (\Cref{prop:vander}).
The bounds are also loose.
Across 5,000 random populations of nine items the left side reached at most $0.271$ of the right side of Equation~\eqref{eq:bounded} and $0.018$ of the variance bound.

\paragraph{Check.}
On 20,000 random pairs of distributions ($K=6$), Pinsker's inequality, the width bound $|\E_pf-\E_qf|\le\Delta_f\TVD(p,q)$ and the variance bound held with smallest slacks $0.0042$, $0.0184$ and $0.7561$, all nonnegative.
The population bounds held in all 5,000 random populations.

\subsection{All moments are the distribution}
\label{app:moments-vander}

\begin{proposition}[The first $K-1$ moments determine the distribution]
\label{prop:vander}
On $K$ rating values the linear map $p\mapsto(\E_pR^j)_{j=0}^{K-1}$ is injective.
Its matrix $V=(r^j)_{r,j=0}^{K-1}$ has determinant $\prod_{k=1}^{K-1}k!$, which equals $34{,}560$ for $K=6$.
Matching the moments $j=0,\dots,s$ leaves an affine set of dimension $K-1-s$; matching the mean alone ($s=1$) leaves $K-2$ free directions.
\end{proposition}

\begin{proof}
For the nodes $0,\dots,K-1$ the Vandermonde determinant is $\prod_{0\le i<j\le K-1}(j-i)=\prod_{j=1}^{K-1}j!\ne0$.
The first $s+1$ columns of $V$ have rank $s+1$, so $s+1$ moment equations leave a solution set of dimension $K-(s+1)$.
\end{proof}

Three consequences follow.
First, ``CE constrains every moment'' and ``CE identifies the distribution'' are one statement.
Second, a mean-only objective leaves $K-2$ shape directions per item, which is the count $n(K-2)+2$ of Proposition~\ref{prop:dimension}.
Third, the moment parameterization is ill-conditioned: $V$ has 2-norm condition number $5.77\times10^{4}$ for $K=6$.
That is one practical reason to fit probabilities directly with a proper loss instead of matching moments one at a time.

\paragraph{Check.}
The exact integer determinant of $V$ for $K=6$ is $34{,}560$, and solving for $p$ from its exact moments recovered 200 random distributions to within $10^{-12}$.
Matching the mean alone leaves 4 free directions, matching the moments up to $K-2$ leaves 1, and matching all $K-1$ leaves 0.

\subsection{Cross-entropy matches moments; the gradient of every objective}
\label{app:moments-grad}

\begin{proposition}[CE stationarity is moment matching]
\label{prop:expfam}
Let $q_\theta(r)\propto\exp(\theta^\top T(r))$ for a feature map $T:\{0,\dots,K-1\}\to\mathbb R^d$.
For one distribution $p$,
\[
\nabla_\theta\,\E_p[-\log q_\theta(R)]=\E_{q_\theta}T(R)-\E_pT(R),
\qquad
\nabla_\theta^2\,\E_p[-\log q_\theta(R)]=\Cov_{q_\theta}\bigl(T(R)\bigr)\succeq0 .
\]
The CE objective is therefore convex in $\theta$, and a point is stationary exactly when the forecast and the human distribution have the same expected features.
\end{proposition}

\begin{proof}
Write $-\log q_\theta(r)=A(\theta)-\theta^\top T(r)$ with the log-partition function $A(\theta)=\log\sum_s\exp(\theta^\top T(s))$.
Its gradient is $\E_{q_\theta}T$ and its Hessian is the covariance of $T$ under $q_\theta$.
\end{proof}

With $T(r)=\bigl((r/(K-1))^j\bigr)_{j=1}^{K-1}$ the family contains every positive distribution, because the matrix of the monomials $1,\dots,(r/(K-1))^{K-1}$ is a Vandermonde matrix and $\log q$ can therefore be any vector up to a constant.
Stationarity then asks for $K-1$ moment equations.
The calibration head of Equation~\eqref{eq:head} is the same construction with $T=(u_k,u_k^2)$ at fixed temperature $t$.
For one observed rating $y$ its intercept and curvature gradients are $\E_qu-u(y)$ and $\E_qu^2-u(y)^2$, so the head matches the first two moments of the centered category index to the rating.

\begin{proposition}[Logit gradients of the objectives]
\label{prop:grad}
Let $q=\mathrm{softmax}(u)$ for one item with mean $m$ and let $y$ be an observed rating.
\Cref{tab:losses} lists the gradients with respect to the logit $u(r)$.
For $n$ items the Pearson gradient is $\partial\rho/\partial u_i(r)=g_i\,q_i(r)(r-m_i)$ with $g_i=(\zeta_i-\rho z_i)/(n\,s_m)$, where $z_i=(m_i-\bar m)/s_m$ and $\zeta_i=(t_i-\bar t)/s_t$ are the standardized forecast and human means.
\end{proposition}

\begin{table}[h]
\caption{\textbf{MSE and Pearson move an item's logits only along the direction that changes its mean; CE, Brier and RPS gradients span all $K-1$ shape directions.}
Notation for one item and one observed rating $y$: $q=\mathrm{softmax}(u)$, $F_k=\sum_{s\le k}q(s)$, $e_y(r)=\Ind\{r=y\}$ and $a_r=\tfrac{2}{K-1}\sum_{k=r}^{K-2}(F_k-\Ind\{y\le k\})$.
The population optimum is the forecast that minimizes the expected loss (maximizes $\rho$ for Pearson); $\mu$ is the true conditional mean.}
\label{tab:losses}
\centering
\small
\begin{tabular}{@{}llll@{}}
\toprule
Objective & Loss & Gradient with respect to $u(r)$ & Population optimum \\
\midrule
CE & $-\log q(y)$ & $q(r)-e_y(r)$ & $q=p$ \\
Brier & $\sum_s\bigl(q(s)-e_y(s)\bigr)^2$ & $2q(r)\bigl[(q(r)-e_y(r))-\sum_sq(s)(q(s)-e_y(s))\bigr]$ & $q=p$ \\
RPS & $\tfrac1{K-1}\sum_k(F_k-\Ind\{y\le k\})^2$ & $q(r)\bigl(a_r-\sum_sq(s)a_s\bigr)$ & $q=p$ \\
MSE & $(m-y)^2$ & $2(m-y)\,q(r)(r-m)$ & $m=\mu$ \\
Pearson & $-\rho(m,t)$ & $-g_i\,q_i(r)(r-m_i)$ & $m=a+b\mu$, $b>0$ \\
\bottomrule
\end{tabular}
\end{table}

The MSE and Pearson gradients are scalar multiples of $(q(r)(r-m))_r$, the gradient of the item's mean itself.
A gradient step on either loss therefore changes an item's logits only along the one direction that changes its mean.
The CE gradient $q-e_y$ ranges over the whole $(K-1)$-dimensional space of logit directions, and averaged over the ratings of one item it equals $q-\hat p$, the gap to the empirical rating histogram.
These are statements about logits.
A network shares parameters across items and categories, its Jacobian mixes these directions, and a correlation gradient can therefore move shape indirectly; Sections~\ref{sec:controlled}--\ref{sec:scaling} measure whether it does.

\paragraph{Check.}
Finite differences reproduce the CE, Brier, MSE and RPS logit gradients (200 random instances each) and the Pearson gradient (100 random instances of six items) to within $10^{-8}$.
The exponential-family gradient and the head gradients agree to the same tolerance.
Polynomial features of degree $K-1$ reproduce 200 random positive distributions to within $10^{-12}$.

\subsection{Pearson is a squared loss on standardized scores}
\label{app:moments-squared}

\begin{proposition}[Pearson is squared error with level and scale removed]
\label{prop:squared}
For forecast means $m_i$ and human means $t_i$ with standard deviations $s_m,s_t$, correlation $\rho$ and standardized values $z_i,\zeta_i$,
\[
\frac1n\sum_i(z_i-\zeta_i)^2=2(1-\rho),
\qquad
\frac1n\sum_i(m_i-t_i)^2=(\bar m-\bar t)^2+(s_m-s_t)^2+2s_ms_t(1-\rho).
\]
\end{proposition}

\begin{proof}
Expand the squares and use $\frac1n\sum_iz_i^2=\frac1n\sum_i\zeta_i^2=1$ and $\frac1n\sum_iz_i\zeta_i=\rho$.
For the second identity, $\frac1n\sum_i(m_i-t_i)^2=(\bar m-\bar t)^2+s_m^2+s_t^2-2\rho s_ms_t$, and $s_m^2+s_t^2-2\rho s_ms_t=(s_m-s_t)^2+2s_ms_t(1-\rho)$.
\end{proof}

Maximizing Pearson correlation therefore minimizes a squared error after the level and the scale of the forecast are chosen freely.
It identifies what squared error identifies, the conditional mean (Proposition~\ref{prop:losses}), up to a positive affine map.

\paragraph{Check.}
Both identities held to within $10^{-13}$ on 200 random pairs of vectors.

\subsection{A finite-sample example with an exact threshold}
\label{app:moments-threshold}

Each item has a rating $R=2+D$ with $D\in\{-1,0,1\}$ and probabilities $\bigl(\tfrac{t-m}{2},\,1-t,\,\tfrac{t+m}{2}\bigr)$.
Here $m$ is a known mean offset, and $t\in\{t_L,t_H\}$ is an unknown spread with $0<t_L<t_H<1$ and $|m|<t_L$.
Both spreads give the mean $2+m$, so a rule that sees only means cannot tell them apart.
We compare two rules that must pick the spread of an item from one rating: the mean-only rule, with its tie broken toward $t_L$, and CE (maximum likelihood), which picks the spread with the larger likelihood.
A fraction $\pi$ of the items has spread $t_H$.

\begin{proposition}[Break-even prior for one rating]
\label{prop:threshold}
The CE rule picks $t_H$ exactly when $D\ne0$.
It errs with probability $e_L=t_L$ on low-spread items and $e_H=1-t_H$ on high-spread items, so its expected error is $(1-\pi)t_L+\pi(1-t_H)$.
The mean-only rule errs with probability $\pi$.
CE has the lower expected error exactly when $\pi>t_L/(t_L+t_H)$.
\end{proposition}

\begin{proof}
For $D=\pm1$ we have $(t_H\mp m)/2>(t_L\mp m)/2$, so the high spread has the larger likelihood, and for $D=0$ we have $1-t_H<1-t_L$.
The CE error $(1-\pi)t_L+\pi(1-t_H)$ is below $\pi$ if and only if $t_L<\pi(t_L+t_H)$.
\end{proof}

For $t_L=0.4$ and $t_H=0.8$ the threshold is $1/3$, with $e_L=0.4$ and $e_H=0.2$.
Wrongly picking the spread costs a squared variance error of $(t_H-t_L)^2=0.16$.
At $\pi=1/3$ both rules lose $0.0533$ in expectation.
At $\pi=0.3$ the mean-only rule loses $0.0480$ against $0.0544$ for CE, and at $\pi=0.4$ CE loses $0.0512$ against $0.0640$ for the mean-only rule.
This is the break-even logic of Section~\ref{sec:controlled} at the smallest scale: with one rating and a prior that strongly favors one shape, ignoring shape is competitive, and once at least a third of the items have the high spread, CE wins.
The example is exact for one rating only.

\paragraph{Check.}
Exact rational arithmetic with five offsets ($m=-0.3,-0.1,0,0.2,0.3$) returns the same $e_L$, $e_H$ and the threshold $1/3=t_L/(t_L+t_H)$ each time.

\subsection{What the derivation shows and does not show}
\label{app:moments-limits}

The derivation concerns population optima and fixed forecasts.
It shows that correlation constrains one conditional moment and that CE constrains all of them, with explicit bounds, gradients and a finite-sample example.
It does not show how a trained network behaves.
Parameters shared between the mean and the shape let a correlation gradient move the shape indirectly, and the bounds are loose in the random instances checked (largest observed ratios $0.271$ and $0.018$).
The random instances are Dirichlet draws, so the checks test identities and inequalities and do not describe real ratings.
Sections~\ref{sec:controlled}--\ref{sec:scaling} and \appref{app:controlled} measure fitted judges on controlled, frozen and trained models.
